\documentclass[10.5pt]{article}
% Shared style for the Airspace Monitor reports
\usepackage[a4paper,margin=1.9cm,top=2.25cm,bottom=2.05cm,headheight=18pt]{geometry}
\usepackage{fontspec}
\usepackage{xcolor}
\usepackage{graphicx}
\usepackage{tikz}
\usetikzlibrary{positioning,calc,shapes.geometric,backgrounds,arrows.meta}
\usepackage[most]{tcolorbox}
\usepackage{booktabs}
\usepackage{amssymb}
\usepackage{tabularx}
\usepackage{longtable}
\usepackage{array}
\usepackage{colortbl}
\usepackage{ragged2e}
\usepackage{enumitem}
\usepackage{fancyhdr}
\usepackage{fontawesome5}
\usepackage{hyperref}
\usepackage{microtype}
\usepackage{lastpage}

\definecolor{ink}{HTML}{0B1020}
\definecolor{panel}{HTML}{111A2E}
\definecolor{border}{HTML}{22304F}
\definecolor{textlight}{HTML}{E7ECF6}
\definecolor{muted}{HTML}{9AA8C2}
\definecolor{accent}{HTML}{4CC3FF}
\definecolor{violet}{HTML}{B392FF}
\definecolor{okgreen}{HTML}{2FAE78}
\definecolor{okgreenlight}{HTML}{E6F7F0}
\definecolor{warn}{HTML}{C77A00}
\definecolor{warnlight}{HTML}{FFF6E5}
\definecolor{alert}{HTML}{C94646}
\definecolor{alertlight}{HTML}{FDEEEE}
\definecolor{paper}{HTML}{F7F9FD}
\definecolor{papertext}{HTML}{1B2438}
\definecolor{papermuted}{HTML}{5B6B85}
\definecolor{lightborder}{HTML}{DCE4F2}
\definecolor{softblue}{HTML}{EAF6FF}
\definecolor{softviolet}{HTML}{F1ECFF}
\definecolor{softgrey}{HTML}{F0F3F9}

\setmainfont{Lato}
\setsansfont{Lato}
\setmonofont{Consolas}[Scale=0.92]
\renewcommand{\familydefault}{\sfdefault}
\setlength{\parindent}{0pt}
\setlength{\parskip}{5pt}
\RaggedRight
\renewcommand{\arraystretch}{1.32}
\setlength{\tabcolsep}{5pt}
\setlength{\emergencystretch}{1em}
\hypersetup{hidelinks}

\newcommand{\sechead}[1]{%
  \vspace{7pt}%
  {\color{papertext}\fontsize{17}{20}\selectfont\bfseries #1}\par
  \vspace{-2pt}{\color{accent}\rule{\linewidth}{1.3pt}}\par
  \vspace{5pt}%
}
\newcommand{\subhead}[1]{%
  \vspace{5pt}{\color{papertext}\fontsize{11.5}{14}\selectfont\bfseries #1}\par
  \vspace{1pt}%
}
\newcommand{\code}[1]{\texttt{\detokenize{#1}}}
\newcommand{\docsrc}[1]{{\scriptsize\color{papermuted}\texttt{\detokenize{#1}}}}

\newcommand{\pill}[2]{%
  \begin{tikzpicture}[baseline=(x.base)]%
    \node[fill=#1,text=white,rounded corners=2.4pt,inner sep=2.6pt,font=\scriptsize\bfseries] (x) {#2};%
  \end{tikzpicture}%
}
\newcommand{\dotok}{\textcolor{okgreen}{\faCircle}}
\newcommand{\dotwarn}{\textcolor{warn}{\faCircle}}
\newcommand{\dotbad}{\textcolor{alert}{\faCircle}}

\newtcolorbox{card}[1][]{enhanced,breakable,colback=white,colframe=lightborder,boxrule=0.6pt,arc=2.2mm,left=4mm,right=4mm,top=3mm,bottom=3mm,#1}
\newtcolorbox{accentcard}[1][]{enhanced,breakable,colback=softblue,colframe=accent,boxrule=0.8pt,arc=2.2mm,left=4mm,right=4mm,top=3mm,bottom=3mm,#1}
\newtcolorbox{warncard}[1][]{enhanced,breakable,colback=warnlight,colframe=warn,boxrule=0.8pt,arc=2.2mm,left=4mm,right=4mm,top=3mm,bottom=3mm,#1}
\newtcolorbox{badcard}[1][]{enhanced,breakable,colback=alertlight,colframe=alert,boxrule=0.8pt,arc=2.2mm,left=4mm,right=4mm,top=3mm,bottom=3mm,#1}
\newtcolorbox{violetcard}[1][]{enhanced,breakable,colback=softviolet,colframe=violet,boxrule=0.8pt,arc=2.2mm,left=4mm,right=4mm,top=3mm,bottom=3mm,#1}
\newtcolorbox{softcard}[1][]{enhanced,breakable,colback=softgrey,colframe=lightborder,boxrule=0.6pt,arc=2.2mm,left=4mm,right=4mm,top=3mm,bottom=3mm,#1}

\pagestyle{fancy}
\fancyhf{}
\fancyhead[L]{\scriptsize\color{papermuted}Airspace Monitor · Documentation Review}
\fancyhead[R]{\scriptsize\color{papermuted}30 October 2026}
\fancyfoot[L]{\scriptsize\color{papermuted}Level 1 · pre-implementation}
\fancyfoot[R]{\scriptsize\color{papermuted}\thepage\ / \pageref{LastPage}}
\renewcommand{\headrulewidth}{0.4pt}
\renewcommand{\footrulewidth}{0.4pt}
\renewcommand{\headrule}{\hbox to\headwidth{\color{lightborder}\leaders\hrule height \headrulewidth\hfill}}
\renewcommand{\footrule}{\hbox to\headwidth{\color{lightborder}\leaders\hrule height \footrulewidth\hfill}}


\begin{document}
\hypersetup{pdftitle={Airspace Monitor - Documentation Review},pdfauthor={Alek Hanusic},pdfsubject={Feasibility review of the Level 1 documentation}}

% ============================================================
% COVER
% ============================================================
\pagecolor{ink}
\color{textlight}
\thispagestyle{empty}

\begin{tikzpicture}[remember picture,overlay]
  % subtle grid
  \foreach \x in {0,...,14} {
    \draw[white,opacity=0.035,line width=0.4pt]
      ($(current page.south west)+(\x*1.55cm,0)$) -- ++(0,29.7cm);
  }
  \foreach \y in {0,...,19} {
    \draw[white,opacity=0.035,line width=0.4pt]
      ($(current page.south west)+(0,\y*1.55cm)$) -- ++(21cm,0);
  }
  % accent orbit
  \draw[accent,opacity=0.35,line width=0.8pt]
    ($(current page.north east)+(-5.2cm,-4.0cm)$) arc[start angle=205,end angle=515,radius=2.35cm];
  \draw[accent,opacity=0.16,line width=0.8pt]
    ($(current page.north east)+(-5.2cm,-4.0cm)$) arc[start angle=205,end angle=515,radius=3.15cm];
  \fill[accent] ($(current page.north east)+(-3.0cm,-5.4cm)$) circle (2.0pt);
  \fill[violet] ($(current page.north east)+(-3.85cm,-2.55cm)$) circle (1.6pt);
  \node[circle,draw=accent,fill=accent!12,minimum size=1.15cm,line width=0.8pt]
    at ($(current page.north east)+(-5.2cm,-4.0cm)$) {\color{accent}\Large\faPlane};
  % bottom accent rule
  \draw[accent,opacity=0.8,line width=1.2pt]
    ($(current page.south west)+(1.9cm,2.0cm)$) -- ++(5.2cm,0);
\end{tikzpicture}

\vspace*{2.2cm}
{\color{accent}\fontsize{9}{11}\selectfont\bfseries DOCUMENTATION REVIEW \quad·\quad LEVEL 1 \quad·\quad 30 OCTOBER 2026}

\vspace{0.55cm}
{\fontsize{40}{44}\selectfont\bfseries Airspace Monitor}

\vspace{0.18cm}
{\fontsize{20}{25}\selectfont Feasibility, risks and the changes\\ I would make before M1}

\vspace{0.9cm}
{\color{muted}\fontsize{11.5}{16}\selectfont
A review of the 11 system documents, four ADRs, the README and the original intent PDF.\par}

\vspace{0.35cm}
{\color{muted}\fontsize{11.5}{16}\selectfont
Prepared for Alek Hanušić \& Dániel Csorba.\par}

\vfill
{\color{muted}\fontsize{9}{12}\selectfont
No source documents were modified. External facts were checked on 30 October 2026.\par
\vspace{2pt}
{\color{accent}\rule{1.2cm}{1pt}}\quad
\textbf{Status:} buildable; public launch is conditional on the OpenSky gate.
}

\newpage
\pagecolor{paper}
\color{papertext}

% ============================================================
% EXECUTIVE VERDICT
% ============================================================
\sechead{Executive verdict}

\begin{warncard}
{\large\bfseries Buildable. Shippable as a public demo only after one licensing gate is cleared.}

\vspace{3pt}
The architecture, data model, pipelines, ML plan and LLM guardrails are strong enough to start building now. The remaining risk is not the technology. It is \textbf{permission} (OpenSky's terms for public display and for two accounts), one \textbf{factual data error} (the Zenodo record), and the \textbf{calendar} (8 weeks for a platform plus three ML products plus an LLM product).
\end{warncard}

\vspace{5pt}
\begin{tabularx}{\linewidth}{@{}l c c X@{}}
\toprule
\textbf{Area} & \textbf{Rating} & \textbf{Score} & \textbf{Assessment} \\
\midrule
Architecture \& hosting & \dotok & 9/10 & Oracle Always Free with a Contabo fallback; multi-arch images; Tailscale/Cloudflare access is realistic. Capacity is the only wildcard. \\
Data model \& pipelines & \dotok & 9/10 & Natural keys, dedupe keys, SI units, explicit job windows, CAGG and retention are well thought through. A few edge cases remain. \\
ML & \dotwarn & 7/10 & ETA is plausible but B1 is a hard baseline; anomaly has no ground truth; route lookup is the right core choice. \\
LLM & \dotok & 8/10 & The SQL guardrails and briefing number check are the strongest part of the design. Eval determinism needs care. \\
Frontend & \dotok & 8/10 & Clear visual identity and sensible layer choices. CARTO key, deck.gl dashed strokes and null-age handling are open details. \\
Ops \& CI & \dotok & 8/10 & Backups, restore procedure, external Grafana sharing and deploy-over-Tailscale are credible. Restore drill and event catch-up are the gaps. \\
Data licensing & \dotbad & 5/10 & OpenSky public redistribution and two-account use need written confirmation. The Zenodo claim is currently a mix of two records. \\
8-week schedule & \dotwarn & 6/10 & Feasible for the core only. Stretch items must stay out until M0--M5 are done. \\
\bottomrule
\end{tabularx}

\vspace{4pt}
\begin{accentcard}
{\bfseries The five actions I would take first}
\begin{enumerate}[leftmargin=1.25em,itemsep=2pt,topsep=2pt]
  \item \textbf{Clear the OpenSky gate in M0.} Ask in writing about public display, a public read-only API, and two personal accounts for one project. Record the answer. If it is no, fall back to one account at 90 s and a private/aggregate-only public demo.
  \item \textbf{Correct the Zenodo claim.} The CC BY 4.0 record (5815448) is 2019--2021 and anonymised. The 2019--2022 record with \code{icao24}/\code{registration} is under OpenSky's licence, not CC BY. Pick one and fix the README and \code{02-data-sources.md}.
  \item \textbf{Fix four small correctness bugs:} \code{ref.runways} empty identifiers, \code{next_expected_at} when one collector is down, the watermark when both collectors are silent, and stream-worker event catch-up.
  \item \textbf{Add the CARTO key} and close the two day-one implementation gaps: gzip request decompression and CORS preview-origin handling.
  \item \textbf{Protect the core.} No fourth ML model and no stretch before M5. ETA beating B1 is the one success criterion that may not be met; plan the negative-result model card now.
\end{enumerate}
\end{accentcard}

\vspace{4pt}
\begin{softcard}
{\bfseries Fixed since the first review.} The current docs already resolve most of the earlier blockers: PostGIS is provided by the pinned \code{timescaledb-ha} image; the density calculation is restricted to global OpenSky snapshots; events have a \code{dedupe_key}; flights have a stable \code{flight_id}; MLflow is private; deploy uses Tailscale; the frontend budget is split into a small shell and a lazy map bundle; and compressed-chunk inserts are guarded by a CI test. That is a meaningful improvement.
\end{softcard}

% ============================================================
% SCOPE AND METHOD
% ============================================================
\sechead{Scope, method and what ``possible'' means}

\begin{minipage}[t]{0.49\linewidth}
\begin{card}
{\bfseries What this review covers}
\begin{itemize}[leftmargin=1.1em,itemsep=1.5pt,topsep=2pt]
  \item \code{docs/00-overview.md} through \code{10-roadmap.md}
  \item ADRs \code{0001}--\code{0004}
  \item \code{README.md}, \code{AGENTS.md}
  \item The original \code{CloudAI-project-intent.pdf}
\end{itemize}
I treated the current Level 1 scope as fixed: OpenSky only, no own ADS-B receiver, two home collectors, Oracle/Contabo hosting, public web app, public Grafana dashboards, private MLflow.
\end{card}
\end{minipage}\hfill
\begin{minipage}[t]{0.49\linewidth}
\begin{card}
{\bfseries The four tests I applied}
\begin{itemize}[leftmargin=1.1em,itemsep=1.5pt,topsep=2pt]
  \item \textbf{Technical:} can it be built with the chosen tools and data?
  \item \textbf{Data:} are the labels, coverage and volume sufficient?
  \item \textbf{Legal:} may the data be stored and shown publicly?
  \item \textbf{Operational:} can two people run it for 8 weeks and keep it alive?
\end{itemize}
A feature is ``possible'' only if all four hold. A green rating means the current plan passes; amber means it needs a change or a measurement; red means it is not true as written or needs a decision.
\end{card}
\end{minipage}

\vspace{5pt}
\begin{softcard}
{\bfseries The idea is coherent.} A live map is the hook; the data platform is the engineering story; ETA and unusual-traffic are the ML story; the guarded text-to-SQL chat is the AI story. The non-goals are sensible (no Kafka/Flink, no schedules, no user accounts). The two weakest points are the public-data permission and the amount of work in 8 weeks, not the architecture.
\end{softcard}

\vspace{3pt}
\begin{minipage}[t]{0.32\linewidth}
\pill{okgreen}{GREEN}\\[2pt]
{\small Possible with the current plan.}
\end{minipage}\hfill
\begin{minipage}[t]{0.32\linewidth}
\pill{warn}{AMBER}\\[2pt]
{\small Possible, but measure or change something first.}
\end{minipage}\hfill
\begin{minipage}[t]{0.32\linewidth}
\pill{alert}{RED}\\[2pt]
{\small Not true as written, or decision required.}
\end{minipage}

\vspace{4pt}
{\small\color{papermuted}
Assumptions: two-person team; 8 weeks starting 27 October 2026; non-commercial MSc project; no code exists yet; external claims were checked against the providers' documentation on 30 October 2026. This is a design review, not legal advice.
}

% ============================================================
% WHAT IS SOLID
% ============================================================
\sechead{What is solid}

\begin{card}
{\bfseries Architecture and hosting.} The ``home collectors push to the cloud'' decision is still the right answer to OpenSky's data-centre blocking, and the ADR now explains why a Tailscale exit node is not the same thing. Oracle Always Free (2 OCPU, 12 GB, 200 GB) is the right cost/scale trade-off, and the Contabo fallback is concrete. Cloudflare for DNS/TLS/rate limits, Tailscale for admin access, no public database port, and multi-arch images are all sensible. The migration rule ``add first, remove later'' is exactly what a deploy-before-code pipeline needs.
\end{card}

\vspace{3pt}
\begin{card}
{\bfseries Data model and pipelines.} This is the strongest part of the docs. Bronze/silver/gold is not decorative: \code{raw} keeps the OpenSky payload normalized but otherwise intact; \code{core} has stable flights and deduplicated events; \code{agg} has a CAGG and a coverage view; \code{ml} and \code{llm} are separate. Natural keys, \code{dedupe_key}, \code{flight_id} that never changes, explicit job windows, and a manual \code{flights_rebuild}/\code{events_rebuild} path are all there. The spool-first collector, R2 archive, catch-up path and DQ checks make the ``nothing is lost, only delayed'' claim mostly true for data that was actually fetched.
\end{card}

\vspace{3pt}
\begin{card}
{\bfseries ML.} The scope is honest: route lookup is core, the re-ranker is stretch; anomaly v1 is statistical; ETA is the only trained model in the core. The ETA design has a real label, a real baseline (B0 distance/speed, B1 distance-band median), time-based split, quantile outputs and coverage metrics. The anomaly design has synthetic injection, persistence and a real-events list because there is no ground truth. Model cards and \code{ops.model_runs} make the results public without exposing MLflow. This is the right level of ambition for 8 weeks.
\end{card}

\vspace{3pt}
\begin{card}
{\bfseries LLM.} The guardrails are the differentiator: a read-only role limited to \code{llm.v_*}, a per-transaction 5 s timeout, a sqlglot allow-list, a 3-call loop cap, a 200-row cap, and the SQL always shown. The briefing is built from a facts JSON and has a number check that rejects hallucinated values. The golden set plus execution accuracy in CI is the right evaluation. This is much stronger than a generic ``chat with your database'' demo.
\end{card}

\vspace{3pt}
\begin{card}
{\bfseries Frontend and ops.} The visual spec is distinctive (dark map, altitude ramp, H3 anomalies, event pulses) and the interaction model is clear. SSE \code{snapshot_ready} plus a cacheable \code{?t=} fetch is the right choice for a 45--90 s feed. The split budget (small shell, lazy map bundle) is realistic. On the ops side, R2 backups, 14-day retention, a documented restore order and externally shared Grafana dashboards are all credible. The public/private boundary is right: MLflow and Dagster stay on the tailnet.
\end{card}

% ============================================================
% GATES
% ============================================================
\sechead{The gates: what is not possible as written}

\begin{badcard}
{\large\bfseries Gate 1 · OpenSky public display and two accounts}

\vspace{3pt}
The docs now correctly say that two accounts for one project must be confirmed against OpenSky's terms. That is necessary but not sufficient. There is a second question the docs do not yet answer: \textbf{may the project publish a public map and a public read-only API that serves OpenSky-derived state vectors?}

The API documentation describes research and non-commercial use, but the data licence is stricter: it limits use to non-profit research/education (plus internal commercial testing and government) and restricts distribution/disclosure of the dataset to collaborators. A public map may qualify as a cited publication; a public \code{/live/snapshot} and \texttt{/aircraft/\{icao24\}/track} serving raw-ish positions is much closer to redistribution. Attribution alone does not resolve this.
\end{badcard}

\vspace{4pt}
\begin{tabularx}{\linewidth}{@{}p{5.0cm} X@{}}
\toprule
\textbf{Question to put to OpenSky in M0} & \textbf{Why it matters / fallback} \\
\midrule
May a non-commercial MSc portfolio display live OpenSky states on a public map with attribution? & If yes, keep the map. If no, the public URL becomes a private demo or a recorded video. \\
May it expose a public read-only \code{/live/snapshot} and \texttt{/aircraft/\{id\}/track} API, or is that redistribution? & If no, keep the raw endpoints behind auth and publish only aggregate/derived views. The map frontend can still use an authenticated endpoint. \\
May two teammates use two personal accounts for one project's redundant collectors? & If no, run one collector at 90 s (the map still moves smoothly via dead reckoning). If yes, keep the 45 s offset and document the permission. \\
\bottomrule
\end{tabularx}

\vspace{4pt}
{\small\color{papermuted}
Action: add this to M0's ``done when'' and keep the written reply in the repo (not the public README). Do not build the public API on the assumption that attribution is enough.}

\vspace{6pt}
\begin{badcard}
{\large\bfseries Gate 2 · The Zenodo record is two different datasets}

\vspace{3pt}
The docs and README currently combine claims that do not belong to one record. This is a factual error, not a judgement call.
\end{badcard}

\vspace{2pt}
\begin{tabularx}{\linewidth}{@{}p{2.5cm} p{4.1cm} p{4.1cm} X@{}}
\toprule
\textbf{Record} & \textbf{What it contains} & \textbf{Licence / coverage} & \textbf{Use it for} \\
\midrule
Zenodo \textbf{5815448} & Monthly \code{flightlist_*.csv.gz}; anonymised aircraft (\code{aircraft_uid}); no \code{icao24} or \code{registration} & \textbf{CC BY 4.0}; January 2019 -- December 2021 & Route priors: only \code{callsign}, \code{origin}, \code{destination} are needed. Safest licence. \\
Zenodo \textbf{3737101} / latest version \textbf{7923702} & Monthly flight lists with \code{icao24} and \code{registration} & \textbf{OpenSky data licence}, not CC BY; January 2019 -- December 2022 & Only if aircraft identity or 2022 routes are required, and only under OpenSky's terms. \\
\bottomrule
\end{tabularx}

\vspace{3pt}
{\small
\textbf{Current claim:} ``OpenSky COVID-19 flight dataset on Zenodo (2019--2022), CC BY 4.0, DOI 10.5281/zenodo.3737101.'' \textbf{Problem:} the CC BY record is 2019--2021 and anonymised; the 2019--2022 record with \code{icao24} is not CC BY. \textbf{Fix:} use 5815448 for route priors, or use 7923702 under the OpenSky licence; correct the README and \code{02-data-sources.md} in the same PR.}

\vspace{6pt}
\begin{warncard}
{\large\bfseries Gate 3 · Other data licences and service keys}

\vspace{2pt}
\begin{itemize}[leftmargin=1.1em,itemsep=2pt,topsep=2pt]
  \item \textbf{OpenSky aircraft database:} non-commercial with attribution, like the API. Record the exact licence text in \code{02} in M3; do not assume it is open data.
  \item \textbf{OpenFlights airlines (ODbL 1.0):} ODbL has share-alike obligations for derivative databases. Attribution in the README may not be sufficient. Verify whether your use is ``insubstantial''; if not, consider a CC0/Wikidata airline list or publish the derived airline table under ODbL.
  \item \textbf{CARTO Dark Matter:} free for non-commercial use up to 5 million tile requests/month, with attribution. Raster tiles now require an API key; vector is recommended to use one. Add the key, a tile budget and the Cloudflare Pages preview origins.
  \item \textbf{OurAirports, Natural Earth:} public domain; no issue. \textbf{H3:} Apache 2.0; no issue.
\end{itemize}
\end{warncard}

\vspace{4pt}
\begin{accentcard}
{\bfseries If OpenSky says no to public redistribution.} You can still ship a strong project: keep the public landing page and demo video, put the live map behind a simple shared password or allowlist for reviewers, and publish only aggregate/derived dashboards publicly. The data platform, models, eval and guardrails remain the portfolio story. Decide this contingency in M0, not in M6.
\end{accentcard}

\vspace{4pt}
\begin{softcard}
{\bfseries Gate status at a glance}
\vspace{2pt}
\begin{itemize}[leftmargin=1.2em,itemsep=2pt,topsep=2pt]
  \item \dotbad\ \textbf{OpenSky public map and API} --- M0 written answer; fallback to a private map or an aggregate-only public view.
  \item \dotwarn\ \textbf{Two accounts for one project} --- M0 written answer; fallback to one collector at 90 s.
  \item \dotbad\ \textbf{Zenodo record} --- choose 5815448 or 7923702; fix DOI, licence and coverage.
  \item \dotwarn\ \textbf{CARTO basemap} --- add API key, tile budget and Pages origins.
  \item \dotwarn\ \textbf{Aircraft DB / OpenFlights} --- record the exact licence; verify ODbL share-alike.
\end{itemize}
\end{softcard}

% ============================================================
% CHANGE LIST
% ============================================================
\newpage
\sechead{Prioritised changes}

{\small
P0 = before M0/M1; P1 = before M2/M3; P2 = before M4/M5 or polish. ``Where'' points to the document that should change. Each item is small enough to land in one PR.
}

\begingroup
\small
\setlength{\tabcolsep}{4pt}
\begin{longtable}{@{}p{0.55cm} >{\raggedright\arraybackslash}p{2.25cm} >{\raggedright\arraybackslash}p{7.25cm} >{\raggedright\arraybackslash}p{4.35cm} >{\raggedright\arraybackslash}p{1.55cm}@{}}
\toprule
\textbf{\#} & \textbf{Where} & \textbf{Change} & \textbf{Why} & \textbf{Priority} \\
\midrule
\endfirsthead
\toprule
\textbf{\#} & \textbf{Where} & \textbf{Change} & \textbf{Why} & \textbf{Priority} \\
\midrule
\endhead
\bottomrule
\endfoot

P0.1 & \code{02}, \code{README}, \code{10} & Resolve OpenSky public-display/API permission and two-account use in M0; record the written answer and add the fallback to the roadmap. & The only true go/no-go. Attribution is not permission to redistribute. & \pill{alert}{P0} \\
P0.2 & \code{02}, \code{README} & Replace the mixed Zenodo claim with one record: 5815448 (CC BY, 2019--2021, anonymised) or 7923702 (OpenSky licence, 2019--2022, with \code{icao24}). & The current DOI/licence/coverage combination is not a real record. & \pill{alert}{P0} \\
P0.3 & \code{03} & \code{ref.runways}: filter rows with empty \code{le_ident}/\code{he_ident} or use a surrogate key; the current PK \code{(airport_icao, le_ident)} fails on OurAirports rows with empty identifiers. & The monthly \code{ref_refresh} will break on real OurAirports data. & \pill{alert}{P0} \\
P0.4 & \code{08} & Add CARTO API key, tile budget and Pages preview origins; note that raster tiles require a key and vector is recommended to use one. & ``Free with attribution'' is still true, but anonymous raster tiles no longer work. & \pill{warn}{P0} \\
P0.5 & \code{05} & Compute \code{next_expected_at} from active collectors/cadence, not a fixed +45 s. & With one collector down the UI promises a 45 s update that will not arrive. & \pill{warn}{P0} \\
P0.6 & \code{04} & Define the watermark when no collector has a successful batch in the last 30 min (fall back to \code{now() - 10 min}). & The current \code{min()} over an empty set can leave flights open forever. & \pill{warn}{P0} \\
P0.7 & \code{04} & Add automatic event catch-up on stream-worker startup (bounded window) or schedule \code{events_rebuild}; do not leave it manual only. & A worker outage longer than 10 min can miss squawks/go-arounds until someone remembers to rebuild. & \pill{warn}{P0} \\
P0.8 & \code{04} & Widen/robustify \code{agg_backfill_refresh} beyond the 65-minute received window, or track unrefreshed batches explicitly. & One failed hourly run can leave late data outside the lookback forever. & \pill{warn}{P0} \\
P0.9 & \code{04}, \code{05} & Add explicit gzip request decompression and enforce the 20 MB limit after decompression. & Starlette's \code{GZipMiddleware} compresses responses; it does not decompress request bodies. & \pill{warn}{P0} \\
P0.10 & \code{05} & Implement CORS preview origins as a regex/allowlist, not a literal \code{https://*.pages.dev} header. & \code{Access-Control-Allow-Origin} does not support wildcard subdomains. & \pill{warn}{P0} \\

P1.1 & \code{04}, \code{06} & Decide what happens to \code{ml.anomalies} when late data changes a bucket: re-score affected buckets, or mark them stale and document it. & \code{anomaly_scoring} never runs for catch-up; the heatmap can disagree with the CAGG. & \pill{warn}{P1} \\
P1.2 & \code{00}, \code{06} & Add a negative-result path for ETA: report MAE vs B0/B1 even if B1 is not beaten, and keep a model card. & B1 is a strong baseline; 4--6 weeks of arrivals may not beat it. The project should not hinge on one metric. & \pill{warn}{P1} \\
P1.3 & \code{06} & Validate H3 resolution 4 against real traffic (or test res 5/6 for a Central Europe layer). & A 1,770 km\textsuperscript{2} cell is coarse for airport-level anomalies; the heatmap may only show overflight patterns. & \pill{warn}{P1} \\
P1.4 & \code{07} & Make \code{now()}-dependent golden questions deterministic: freeze the fixture clock or use fixed reference dates. & Execution accuracy will be flaky if the reference SQL and the model query see different ``now''. & \pill{warn}{P1} \\
P1.5 & \code{08} & Handle \code{pos_age_s = null} in dead reckoning; do not extrapolate unknown-age states. Confirm the dashed H3 outline with a separate line layer if \code{H3HexagonLayer} cannot dash. & Small but visible map bugs; the visual spec depends on them. & \pill{warn}{P1} \\
P1.6 & \code{02}, \code{README} & Record the OpenSky aircraft DB exact licence and verify OpenFlights ODbL share-alike; attribution may not be enough for a public app. & Two data sources are currently ``credit only'' without a licence check. & \pill{warn}{P1} \\
P1.7 & \code{01}, \code{04} & Keep the M1 restore drill as a hard gate; test \code{pg_dump} excluding Timescale chunk data and the replay from R2 on the pinned image. & The backup is the riskiest ops step and the one most likely to be discovered too late. & \pill{warn}{P1} \\
P1.8 & \code{08} & Measure the 150 KB shell and 800 KB lazy map budgets in M2; lazy-load panels and deck.gl layers if needed. & React 19 + TanStack + Tailwind + UI can exceed 150 KB; the map bundle is the bigger risk. & \pill{warn}{P1} \\

P2.1 & \code{README} & Replace ``every plane in the sky'' with ``aircraft visible to OpenSky's volunteer network''; replace ``nothing is lost'' with ``data already fetched is never lost; gaps are recorded''; clarify the \$40 hard cap vs the \$2--8 typical spend; make 45 s conditional. & The opening overclaims coverage and reliability; the details are more honest than the summary. & \pill{okgreen}{P2} \\
P2.2 & \code{10} & Add the licensing gate, CARTO key and restore drill to the M0/M1 ``done when'' lists. & These are currently implicit; the roadmap should make them explicit. & \pill{okgreen}{P2} \\
P2.3 & \code{00} & Add a reference device/connection to the ``<3 s / 60 fps'' criterion, and a fallback for the ETA criterion. & Otherwise ``done'' is untestable on a different laptop. & \pill{okgreen}{P2} \\
P2.4 & \code{06} & Frame route lookup as a deterministic prior, not ML; keep the ML story on ETA + anomaly. & It is a good core feature, but it is not a trained model. & \pill{okgreen}{P2} \\
P2.5 & \code{05} & Make \code{llm.v_live_aircraft} respect the 10-minute live window or label stale answers. & A chat answer from a stale snapshot should say so. & \pill{okgreen}{P2} \\

\end{longtable}
\endgroup

\begin{softcard}
{\bfseries What I would not change.} Keep the single Timescale/PostGIS store, Dagster, the spool-first collector, the read-only chat role, the private MLflow, the Cloudflare/Tailscale boundary and the core/stretch split. Those are the decisions that make the project explainable in an interview. The fixes above are edge cases and permissions, not a redesign.
\end{softcard}

\vspace{3pt}
\begin{softcard}
{\bfseries Suggested sequencing}
\vspace{2pt}
\begin{itemize}[leftmargin=1.2em,itemsep=2pt,topsep=2pt]
  \item \textbf{Week 0} --- P0.1--P0.4: permissions, Zenodo, runways, CARTO.
  \item \textbf{Week 1} --- P0.5--P0.10: API and pipeline correctness before the public URL.
  \item \textbf{M2} --- P1.1--P1.5: event catch-up, ETA fallback, H3 check, eval time, frontend details.
  \item \textbf{M4} --- P1.6--P1.8: remaining licences, restore drill, performance budgets.
  \item \textbf{M5--M6} --- P2 polish; do not start stretch.
\end{itemize}
\end{softcard}

\vspace{3pt}
\begin{softcard}
{\bfseries Questions to settle with the team in week 0}
\vspace{2pt}
\begin{itemize}[leftmargin=1.2em,itemsep=2pt,topsep=2pt]
  \item Who owns Track A and Track B (O1)?
  \item Which project name and domain (O5)?
  \item Is a dedicated low-power collector device approved (O7)?
  \item If two OpenSky accounts are not allowed, which collector keeps polling?
  \item Who writes the ETA model card and the README demo script?
\end{itemize}
\end{softcard}

% ============================================================
% SCOPE AND SCHEDULE
% ============================================================
\newpage
\sechead{Scope, schedule and the cut line}

\begin{tabularx}{\linewidth}{@{}p{1.3cm} X X@{}}
\toprule
\textbf{M} & \textbf{Main risk} & \textbf{Mitigation} \\
\midrule
M0 & OpenSky terms, Oracle capacity, domain, CARTO key & Ask all three vendors in week 0; keep written answers; order the Contabo fallback if A1 is not allocated by end of week 1. \\
M1 & Restore drill, compression migration, backlog drain & Do not merge the compression policy until the backlog is drained; run the restore drill before M2. \\
M2 & Map performance, Cloudflare/SSE, CORS & Measure 12k icons and first load on the reference laptop; test SSE through Cloudflare early. \\
M3 & Flight labels, route coverage, reference loaders & Audit 100 random flights; filter runways; keep the 20-minute boundary rules. \\
M4 & ETA vs B1, anomaly no ground truth & Start collection now; define a negative result; report synthetic recall and events/day. \\
M5 & Chat execution accuracy, cost & Freeze fixture time; keep the 3-call cap and \$40 spend limit; test 30+ questions. \\
M6 & Polish eating the buffer & Freeze features at end of M5; use M6 for README, video and model cards only. \\
\bottomrule
\end{tabularx}

\vspace{4pt}
\begin{minipage}[t]{0.49\linewidth}
\begin{accentcard}
{\bfseries Protect (core)}
\begin{itemize}[leftmargin=1.05em,itemsep=1.5pt,topsep=2pt]
  \item Collector, spool, R2 archive
  \item Ingest idempotency and DQ
  \item Live map + events
  \item Anomaly v1 + heatmap
  \item ETA + arrivals board
  \item Route lookup
  \item Briefing + guarded chat + eval
  \item Grafana ops/models dashboards
\end{itemize}
\end{accentcard}
\end{minipage}\hfill
\begin{minipage}[t]{0.49\linewidth}
\begin{warncard}
{\bfseries Defer (stretch)}
\begin{itemize}[leftmargin=1.05em,itemsep=1.5pt,topsep=2pt]
  \item Destination re-ranker
  \item Anomaly v2 (STL / Isolation Forest)
  \item Holding-pattern events
  \item Self-hosted PMTiles
  \item Any fourth ML model
  \item Level 2 ADS-B receiver (out of scope)
\end{itemize}
\end{warncard}
\end{minipage}

\vspace{4pt}
\begin{warncard}
{\bfseries The one criterion that may fail.} ``ETA p50 MAE beats B1'' is a real research result, not a delivery task. B1 is a distance-band median from the same data; with 4--6 weeks it can be strong. Keep the baseline, report the result honestly, and make the model card useful even if the challenger is not promoted. A well-analysed near-miss is a better portfolio story than a hidden failure.
\end{warncard}

\vspace{2pt}
{\small\color{papermuted}
\textbf{8-week verdict:} core is feasible for two people if this is the main project and collection starts now. The stretch list is not feasible in the same 8 weeks; it should stay out until M5 is green. The schedule risk is concentrated in M1 (platform), M2 (map), and M4 (ETA); M6 is not a buffer for features.
}

\vspace{4pt}
\begin{softcard}
{\bfseries Four rules to keep the scope honest}
\begin{itemize}[leftmargin=1.1em,itemsep=2pt,topsep=2pt]
  \item No stretch until M0--M5 are green.
  \item No public raw OpenSky data until the M0 permission is recorded.
  \item No compression policy until the M1 backlog is drained.
  \item No fourth ML model. ETA and unusual traffic are enough for this project.
\end{itemize}
\end{softcard}

% ============================================================
% DOCUMENT-BY-DOCUMENT
% ============================================================
\newpage
\sechead{Document-by-document notes}

\begingroup
\small
\setlength{\tabcolsep}{4pt}
\begin{longtable}{@{}>{\raggedright\arraybackslash}p{2.9cm} p{1.3cm} >{\raggedright\arraybackslash}p{12.5cm}@{}}
\toprule
\textbf{Document} & \textbf{Verdict} & \textbf{Notes and change} \\
\midrule
\endfirsthead
\toprule
\textbf{Document} & \textbf{Verdict} & \textbf{Notes and change} \\
\midrule
\endhead
\bottomrule
\endfoot
\code{00-overview.md} & \dotwarn & Strong scope, success criteria and known limits. Add the OpenSky public-display gate; make the 45 s cadence conditional on the two-account answer; add a reference device for the map criterion and a negative-result fallback for ETA. \\
\code{01-architecture.md} & \dotok & Oracle/Contabo, Tailscale, Cloudflare and the deploy order are coherent. Keep the restore drill as a gate; verify \code{pg_dump} excludes chunk data on the pinned image. \\
\code{02-data-sources.md} & \dotbad & Good field-level detail and SI METAR. Fix the Zenodo record; add the OpenSky public-redistribution question; record the aircraft DB licence and verify OpenFlights ODbL. \\
\code{03-data-model.md} & \dotok & The strongest document. Natural keys, dedupe, retention and roles are clear. Fix \code{ref.runways}; add a watermark fallback and decide late-data anomaly re-scoring. \\
\code{04-pipelines.md} & \dotok & Spool, catch-up, idempotency and DQ are well designed. Add gzip decompression, stream-worker event catch-up, a wider backfill lookback, and the empty-watermark case. \\
\code{05-api-contract.md} & \dotok & Clear boundary and good cache/SSE design. Fix \code{next_expected_at}; make CORS previews a regex; add request decompression; make the live view freshness explicit. \\
\code{06-ml.md} & \dotwarn & Right scope and good baselines. ETA may not beat B1; H3 res 4 may be coarse; late data can stale the heatmap; route lookup is not ML. \\
\code{07-llm.md} & \dotok & Excellent guardrails. Freeze fixture time for \code{now()} questions; label stale live answers; keep the cost cap. \\
\code{08-frontend.md} & \dotok & Strong visual spec. Add the CARTO key; handle null \code{pos_age_s}; confirm the dashed H3 outline; measure the two budgets. \\
\code{09-conventions.md} & \dotok & Clear tooling and migration rules. Add the gzip/licensing notes; keep the stretch freeze. \\
\code{10-roadmap.md} & \dotok & Good milestone structure. Add the licensing gate, CARTO key and restore drill to the ``done when'' lists; protect the cut line. \\
ADR 0001 & \dotwarn & Correct home-collector decision. Add the public-display permission and the conditional 45 s. \\
ADR 0002 & \dotok & The revised ``Level 2 is out of scope'' is honest and specific. Keep it. \\
ADR 0003 & \dotok & Good single-store decision; the pinned 2.19+ and compressed-chunk test are the right safeguards. \\
ADR 0004 & \dotok & Realistic Oracle choice with a concrete fallback. Capacity is the only open risk. \\
\code{README.md} & \dotwarn & The plain-language story is good, but the opening overclaims ``every plane'' and ``nothing is lost''. Fix the Zenodo claim, the \$40 cap and the conditional 45 s. \\
\code{AGENTS.md} & \dotok & Clear scope, idempotency and migration rules. Consider adding ``do not expose raw OpenSky data publicly until the M0 permission is recorded''. \\
\end{longtable}
\endgroup

% ============================================================
% M0 CHECKLIST AND SOURCES
% ============================================================
\sechead{Closing assessment}

\begin{accentcard}
{\large\bfseries Build it, but clear the gate first.}

\vspace{3pt}
The docs are past the ``is this a good idea?'' stage. The architecture is appropriate, the data model is strong, and the LLM guardrails are unusually good for a student project. The remaining risks are external: OpenSky's permission, one data-source fact, and the ETA result. Resolve the first two in week 0; plan for the third in the model card.

\vspace{3pt}
\textbf{Success at the end of week 8:} a public map that is honest about coverage; one ML result with a baseline and a model card; a chat that shows its SQL and passes the fixed eval; and a stranger can run the stack from the README.
\end{accentcard}

\vspace{5pt}
\newpage
\sechead{M0 go/no-go checklist}

\begin{card}
\begin{itemize}[leftmargin=1.55em,itemsep=3pt,topsep=2pt]
  \item[$\square$] Written OpenSky answer: public map, public API, two accounts.
  \item[$\square$] If two accounts are not allowed: one collector at 90 s; update README/00/01.
  \item[$\square$] If public API is not allowed: private map or aggregate-only public view; update the success criteria.
  \item[$\square$] Zenodo record chosen and DOI/licence/coverage corrected in README/02.
  \item[$\square$] \code{ref.runways} loader filters empty identifiers; a test fixture covers it.
  \item[$\square$] CARTO key created, tile budget set, Pages preview origins allowed.
  \item[$\square$] Oracle A1 allocated or Contabo fallback ordered.
  \item[$\square$] Domain registered; Cloudflare and Tailscale access tested.
  \item[$\square$] Restore drill scheduled for M1, not M6.
\end{itemize}
\end{card}

\vspace{4pt}
\begin{softcard}
{\bfseries Sources checked (30 October 2026).}
\begin{itemize}[leftmargin=1.05em,itemsep=1.5pt,topsep=2pt]
  \item OpenSky REST API (OAuth2, 4,000/8,000 credits, 4 credits per global call): \url{https://openskynetwork.github.io/opensky-api/rest.html}
  \item OpenSky terms of use: \url{https://opensky-network.org/about/terms-of-use}
  \item Zenodo CC BY record: \url{https://zenodo.org/records/5815448}; original/latest record: \url{https://zenodo.org/records/7923702}
  \item CARTO basemaps FAQ (free tier, key requirement): \url{https://docs.carto.com/faqs/carto-basemaps.md}
  \item Grafana externally shared dashboards (OSS support): \url{https://grafana.com/docs/grafana/latest/visualizations/dashboards/share-dashboards-panels/shared-dashboards/}
  \item TimescaleDB HA image (PostGIS included): \url{https://github.com/timescale/timescaledb-docker-ha}
  \item Timescale compressed-chunk \code{ON CONFLICT} issue: \url{https://github.com/timescale/timescaledb/issues/7672}
  \item OurAirports runways data (empty identifiers exist): \url{https://davidmegginson.github.io/ourairports-data/runways.csv}
  \item Claude Haiku 5.5 model page: \url{https://platform.claude.com/docs/en/models/haiku-5-5/overview}
  \item MapLibre bundle size reference: \url{https://bundlephobia.com/package/maplibre-gl@5.0.0}
\end{itemize}
\end{softcard}

\vspace{3pt}
{\scriptsize\color{papermuted}
This review is based on the documents as of 30 October 2026. It does not replace a legal review of the data licences; the OpenSky and ODbL points should be confirmed with the providers. The author did not modify the source documents.
}

\end{document}
